% docs/manual/latex/user/07_outputs.tex - chapter of the LaTeX edition of the book "user_manual" (SIMC++ documentation).
%   Converted once from docs/manual/markdown/user/07_outputs.md by tools/md2latex.py; from then on THIS FILE is the source of the
%   LaTeX edition and may be edited by hand. Re-running tools/md2latex.py --write-sources for this
%   chapter OVERWRITES the hand edits (see docs/README.md).
%   Compile: cmake --build <build> --target pdf, or tools/build_docs_pdf.py (figures, xelatex passes,
%   makeindex with docs/common/latex/simcpp.ist); by hand, see docs/README.md. XeLaTeX, class
%   docs/common/latex/simcpp.cls (packages: fontspec, tikz, xcolor, hyperref, listings, longtable, booktabs,
%   titlesec, titletoc, fancyhdr, makeidx, preview, ...); fonts TeX Gyre Pagella, TeX Gyre Heros, DejaVu Sans
%   Mono. The figures are built first from docs/{manual,translation}/figures/*.tex (preview.sty) into PDFs.
\chapter{\texorpdfstring{Outputs}{Outputs}}\label{user_manual-07u5foutputs}
\simcabstract{A run writes into the output directory (\code{-{}o\ DIR}, default \code{.}), with \code{<run>} = the input file name without its extension
(or \code{-{}-{}run-{}name}\simcidx{Command-line options!run-name@\code{-{}-{}run-{}name}}):}
\begin{simctable}{small}{3}{>{\raggedright\arraybackslash}p{\dimexpr 0.250\linewidth-2\tabcolsep\relax}>{\raggedright\arraybackslash}p{\dimexpr 0.209\linewidth-2\tabcolsep\relax}>{\raggedright\arraybackslash}p{\dimexpr 0.473\linewidth-2\tabcolsep\relax}}
%
\toprule\rowcolor{simchead}\simcth{file} & \simcth{written by (sink)} & \simcth{content} \\\midrule\endfirsthead
\multicolumn{3}{@{}l}{\simccontinued}\\\toprule\rowcolor{simchead}\simcth{file} & \simcth{written by (sink)} & \simcth{content} \\\midrule\endhead
\midrule\multicolumn{3}{r@{}}{\simccontinues}\\\endfoot
\bottomrule\endlastfoot
\code{<run>\allowbreak{}.\allowbreak{}root} & \code{root\_\allowbreak{}tree} & tree \code{h10}, one entry per accepted event; only if \code{Nntu\ >\ 0} \\
\code{<run>\allowbreak{}.\allowbreak{}hist} & \code{fortran\_\allowbreak{}hist} & the run report \\
\code{<run>\allowbreak{}.\allowbreak{}geni} & \code{fortran\_\allowbreak{}hist} & spectrometer stop counters \\
\code{<run>\allowbreak{}.\allowbreak{}gen} & \code{fortran\_\allowbreak{}hist} & diagnostic histograms (text) \\
\code{<run>\allowbreak{}.\allowbreak{}gen.\allowbreak{}root} & \code{gen\_\allowbreak{}hist} & the same histograms as ROOT TH1D \\
\code{<run>\allowbreak{}.\allowbreak{}summary.\allowbreak{}json}, \code{.\allowbreak{}summary.\allowbreak{}txt} & \code{run\_\allowbreak{}summary} & run summary (schema \code{docs/\allowbreak{}manual/\allowbreak{}markdown/\allowbreak{}reference/\allowbreak{}summary.\allowbreak{}schema.\allowbreak{}json}) \\
\end{simctable}
\code{ngen\ =\ 0} writes nothing. The detailed reference of every value is \code{docs/\allowbreak{}manual/\allowbreak{}markdown/\allowbreak{}reference/\allowbreak{}OUTPUT\_\allowbreak{}FORMAT.\allowbreak{}md}.

\section{\texorpdfstring{The event tree \code{h10}}{The event tree h10}}\label{user_manual-07u5foutputs:71-the-event-tree-h10}
\begin{itemize}
\item One entry per \textbf{accepted} event.
\item Every branch is a \code{Float\_\allowbreak{}t} (computed in double precision, stored as float). The branch names and order are those of
the classic SIMC ntuple, so existing analysis scripts work unchanged.
\item The layout depends on the reaction (\code{NtupleColumnsFor}):
\end{itemize}
\begin{simctable}{small}{3}{>{\raggedright\arraybackslash}p{\dimexpr 0.584\linewidth-2\tabcolsep\relax}>{\raggedright\arraybackslash}p{\dimexpr 0.136\linewidth-2\tabcolsep\relax}>{\raggedright\arraybackslash}p{\dimexpr 0.212\linewidth-2\tabcolsep\relax}}
%
\toprule\rowcolor{simchead}\simcth{run} & \simcth{columns} & \simcth{section} \\\midrule\endfirsthead
\multicolumn{3}{@{}l}{\simccontinued}\\\toprule\rowcolor{simchead}\simcth{run} & \simcth{columns} & \simcth{section} \\\midrule\endhead
\midrule\multicolumn{3}{r@{}}{\simccontinues}\\\endfoot
\bottomrule\endlastfoot
H/\allowbreak{}D/\allowbreak{}A(e,\allowbreak{}e'p) and \hyperref[param:doingu5fphsp]{\code{doing\_\allowbreak{}phsp}}\simcidx{Parameters!doing_phsp@\code{doing\_phsp}} & 46 & \hyperref[user_manual-07u5foutputs:713-eep-layout-46-columns]{7.1.3} \\
pion (\hyperref[param:doingu5fpion]{\code{doing\_\allowbreak{}pion}}\simcidx{Parameters!doing_pion@\code{doing\_pion}}), Delta (\hyperref[param:doingu5fdelta]{\code{doing\_\allowbreak{}delta}}\simcidx{Parameters!doing_delta@\code{doing\_delta}}) & 53; π0: 65 & \hyperref[user_manual-07u5foutputs:714-pion-kaon-and-delta-layouts]{7.1.4} \\
kaon (\hyperref[param:doingu5fkaon]{\code{doing\_\allowbreak{}kaon}}\simcidx{Parameters!doing_kaon@\code{doing\_kaon}}) & 55 & \hyperref[user_manual-07u5foutputs:714-pion-kaon-and-delta-layouts]{7.1.4} \\
SIDIS (\hyperref[param:doingu5fsemi]{\code{doing\_\allowbreak{}semi}}\simcidx{Parameters!doing_semi@\code{doing\_semi}}) & 56; π0: 68 & \hyperref[user_manual-07u5foutputs:715-sidis-layout-56-columns]{7.1.5} \\
rho (\hyperref[param:doingu5frho]{\code{doing\_\allowbreak{}rho}}\simcidx{Parameters!doing_rho@\code{doing\_rho}}) & 59 & \hyperref[user_manual-07u5foutputs:716-rho-layout-59-columns]{7.1.6} \\
any of the above except (e,e'p), with \hyperref[param:usingu5ftgtu5ffield]{\code{using\_\allowbreak{}tgt\_\allowbreak{}field}}\simcidx{Parameters!using_tgt_field@\code{using\_tgt\_field}} & + 8 & \hyperref[user_manual-07u5foutputs:717-polarized-target-columns-usingu5ftgtu5ffield]{7.1.7} \\
\end{simctable}
\subsection{\texorpdfstring{Weights and normalisation}{Weights and normalisation}}\label{user_manual-07u5foutputs:711-weights-and-normalisation}
\begin{itemize}
\item \code{Weight} is the event weight computed for each event: the cross section of the event times the generation
weight (radiative weight, jacobians, spectral-function or survival factors, depending on the reaction).
\item The yield of one event for the simulated charge \hyperref[param:EXPERu25charge]{\code{EXPER\%\allowbreak{}charge}}\simcidx{Parameters!EXPER.charge@\code{EXPER\%charge}} is \code{Weight\ ×\ normfac\ /\ nevent}. \code{normfac} is printed in
the \code{.\allowbreak{}hist} and the summary, and \code{nevent} = Ncontribute when ngen > 0, Ntried when ngen < 0.
\code{normfac\ =\ luminosity\ ×\ genvol\ /\ Ntried\ ×\ nevent}, and 1 for \hyperref[param:doingu5fphsp]{\code{doing\_\allowbreak{}phsp}}\simcidx{Parameters!doing_phsp@\code{doing\_phsp}}.
\item The \code{.\allowbreak{}hist} \code{wtcontr} is the sum of these yields over the accepted events that pass the \code{SPedge} window
(pass\_\allowbreak{}cuts): the expected number of counts.
\item \code{sigcc} / \code{siglab} is the cross section of the event (units in the tables below).
\end{itemize}
\subsection{\texorpdfstring{Arm prefixes \code{hs*} and \code{ss*}}{Arm prefixes hs* and ss*}}\label{user_manual-07u5foutputs:712-arm-prefixes-hs-and-ss}
As in results\_\allowbreak{}write.\allowbreak{}f,\allowbreak{} \code{hs*} is the spectrometer on the beam \textbf{right} and \code{ss*} the one on the \textbf{left}:

\begin{itemize}
\item \hyperref[param:electronu5farm]{\code{electron\_\allowbreak{}arm}}\simcidx{Parameters!electron_arm@\code{electron\_arm}} = 1, 3, 7 (HMS, HRS-R, calorimeter HMS side): \code{hs*} = electron, \code{ss*} = hadron;
\item \hyperref[param:electronu5farm]{\code{electron\_\allowbreak{}arm}}\simcidx{Parameters!electron_arm@\code{electron\_arm}} = 2, 4, 5, 8 (SOS, HRS-L, SHMS, calorimeter SHMS side): \code{hs*} = hadron, \code{ss*} = electron.
\end{itemize}
Slopes (\code{*yptar}, \code{*xptar}, \code{*xpfp}, \code{*ypfp}) are in radians (dx/dz); deltas are in \%.

\subsection{\texorpdfstring{(e,e'p) layout (46 columns)}{(e,e'p) layout (46 columns)}}\label{user_manual-07u5foutputs:713-eep-layout-46-columns}
\begin{simctable}{small}{4}{>{\raggedright\arraybackslash}p{\dimexpr 0.073\linewidth-2\tabcolsep\relax}>{\raggedright\arraybackslash}p{\dimexpr 0.247\linewidth-2\tabcolsep\relax}>{\raggedright\arraybackslash}p{\dimexpr 0.443\linewidth-2\tabcolsep\relax}>{\raggedright\arraybackslash}p{\dimexpr 0.148\linewidth-2\tabcolsep\relax}}
\caption{(e,e'p) layout (46 columns)}\\%
\toprule\rowcolor{simchead}\simcth{\#} & \simcth{branches} & \simcth{content} & \simcth{units} \\\midrule\endfirsthead
\multicolumn{4}{@{}l}{\simccontinued}\\\toprule\rowcolor{simchead}\simcth{\#} & \simcth{branches} & \simcth{content} & \simcth{units} \\\midrule\endhead
\midrule\multicolumn{4}{r@{}}{\simccontinues}\\\endfoot
\bottomrule\endlastfoot
1-4 & \code{hsdelta\ hsyptar\ hsxptar\ hsytar} & right arm, reconstructed & \%, rad, rad, cm \\
5-8 & \code{hsxfp\ hsxpfp\ hsyfp\ hsypfp} & right arm focal plane & cm, rad, cm, rad \\
9-12 & \code{hsdeltai\ hsyptari\ hsxptari\ hsytari} & right arm, generated (at the vertex, before the spectrometer) & \%, rad, rad, cm \\
13-24 & \code{ssdelta\ .\allowbreak{}.\allowbreak{}.\ ssytari} & the same for the left arm &  \\
25-29 & \code{q\ nu\ Q2\ W\ epsilon} & reconstructed kinematics & GeV/c, GeV, GeV², GeV, — \\
30-31 & \code{Em\ Pm} & missing energy and momentum & GeV, GeV/c \\
32-33 & \code{thetapq\ phipq} & angles of the proton w.r.t. q & rad \\
34 & \code{corrsing} & e' momentum minus the elastic value & GeV \\
35-37 & \code{Pmx\ Pmy\ Pmz} & missing-momentum components ("Heepcheck" frame) & GeV/c \\
38-40 & \code{PmPar\ PmPer\ PmOop} & missing momentum parallel, perpendicular,\allowbreak{} out of plane & GeV/c \\
41 & \code{fry} & −(raster vertical position) & cm \\
42 & \code{radphot} & radiated photon energy & GeV \\
43 & \code{sigcc} & cross section (µb/sr for H, µb/MeV/sr² for D/A) &  \\
44 & \code{Weight} & event weight &  \\
45-46 & \code{theta\_\allowbreak{}e\ theta\_\allowbreak{}p} & reconstructed scattering angles & rad \\
\end{simctable}
\subsection{\texorpdfstring{Pion, kaon and Delta layouts}{Pion, kaon and Delta layouts}}\label{user_manual-07u5foutputs:714-pion-kaon-and-delta-layouts}
Columns 1-33 as above, then:

\begin{simctable}{small}{4}{>{\raggedright\arraybackslash}p{\dimexpr 0.073\linewidth-2\tabcolsep\relax}>{\raggedright\arraybackslash}p{\dimexpr 0.227\linewidth-2\tabcolsep\relax}>{\raggedright\arraybackslash}p{\dimexpr 0.430\linewidth-2\tabcolsep\relax}>{\raggedright\arraybackslash}p{\dimexpr 0.181\linewidth-2\tabcolsep\relax}}
\caption{Pion, kaon and Delta layouts}\\%
\toprule\rowcolor{simchead}\simcth{\#} & \simcth{branch} & \simcth{content} & \simcth{units} \\\midrule\endfirsthead
\multicolumn{4}{@{}l}{\simccontinued}\\\toprule\rowcolor{simchead}\simcth{\#} & \simcth{branch} & \simcth{content} & \simcth{units} \\\midrule\endhead
\midrule\multicolumn{4}{r@{}}{\simccontinues}\\\endfoot
\bottomrule\endlastfoot
34 & \code{missmass} & missing mass (on the struck nucleon) & GeV \\
35 & \code{mmnuc} & missing mass on the nucleus & GeV \\
36 & \code{phad} & hadron momentum & GeV/c \\
37 & \code{t} & −t (SIMC sign), reconstructed & GeV² \\
38-40 & \code{pmpar\ pmper\ pmoop} & missing-momentum components & GeV/c \\
41 & \code{fry} & −(raster vertical position) & cm \\
42 & \code{radphot} & radiated photon energy & GeV \\
43 & \code{pfermi} & Fermi momentum along q (signed) & GeV/c \\
44 & \code{siglab} & lab cross section & µb/MeV/sr² \\
45 & \code{sigcm} & CM cross section & model units \\
46 & \code{Weight} & event weight &  \\
47 & \code{decdist} & decay distance (cm); for kaons with decay OFF the survival probability &  \\
48 & \code{Mhadron} & mass of the detected particle (μ or π after a decay) & MeV \\
49 & \code{pdotqhat} & Fermi momentum · q̂ & GeV/c \\
50-53 & \code{Q2i\ Wi\ ti\ phipqi} & vertex (generated) Q², W, t, φpq & GeV², GeV, GeV², rad \\
54-65 & \code{xcal\_\allowbreak{}gamma1\ ycal\_\allowbreak{}gamma1\ Egamma1\ Pgamma1x\ Pgamma1y\ Pgamma1z} and the same for photon 2 & π0 runs only: photon positions on the calorimeter (cm; −1e10 = missed, 0 with HMS/SHMS) and lab photon energy and momentum (MeV) &  \\
\end{simctable}
Kaon runs add 54 \code{saghai} (Saghai-model cross section) and 55 \code{factor}
(factorized model, = \code{sigcm}). Delta runs use the 53 pion columns: \code{missmass} is the missing pion mass, \code{phad} the proton
momentum, \code{siglab} and \code{Weight} contain \textbf{no cross section} (\hyperref[user_manual-05u5freactions:55-eepu3c0-in-the-delta-region-doingu5fdelta]{5.5}),
and \code{decdist}, \code{Mhadron} are 0 and Mp.

\subsection{\texorpdfstring{SIDIS layout (56 columns)}{SIDIS layout (56 columns)}}\label{user_manual-07u5foutputs:715-sidis-layout-56-columns}
Columns 1-33 as (e,e'p), then: 34 \code{missmass} (GeV), 35 \code{ppi} (hadron momentum, GeV/c), 36 \code{t} (GeV²), 37 \code{fry}, 38
\code{radphot}, 39 \code{siglab} (µb/\allowbreak{}MeV²/\allowbreak{}sr²),\allowbreak{} 40 \code{sigcent} (always 0), 41 \code{Weight}, 42 \code{decdist} (cm, or the survival probability
with decay OFF), 43 \code{Mhadron} (MeV), 44-45 \code{z\ zi}, 46-47 \code{pt2\ pt2i} (GeV²), 48-49 \code{xbj\ xbji}, 50 \code{thqi} (rad), 51 \code{sighad}
(GeV\textsuperscript{−}²),\allowbreak{} 52 \code{jacobian}, 53 \code{centjac} (always 0), 54 \code{pfermi} (GeV/c), 55 \code{xfermi}, 56 \code{phipqi} (rad). π0 SIDIS appends the 12 photon columns (57-68). The \code{i}
suffix means vertex (generated) values.

\subsection{\texorpdfstring{Rho layout (59 columns)}{Rho layout (59 columns)}}\label{user_manual-07u5foutputs:716-rho-layout-59-columns}
The 56 SIDIS columns with the rho meanings, then 57 \code{Mrho} (Breit-Wigner ρ mass of the event, MeV), 58 \code{Thrho} (decay
angle in the ρ rest frame w.r.t. the lab z axis, rad) and 59 \code{mmnuc} (GeV). The hadron-arm columns describe the detected
decay pion; \code{t} is the pion t; \code{siglab} (µb/MeV/sr) is the ρ cross section; \code{sigcent}, \code{zi}, \code{pt2i}, \code{xfermi} are 0.

\subsection{\texorpdfstring{Polarized-target columns (\code{using\_\allowbreak{}tgt\_\allowbreak{}field})}{Polarized-target columns (using\_\allowbreak{}tgt\_\allowbreak{}field)}}\label{user_manual-07u5foutputs:717-polarized-target-columns-usingu5ftgtu5ffield}
Eight columns follow \code{phipqi} (pion/\allowbreak{}kaon/\allowbreak{}Delta: 54-61; SIDIS/rho: 57-64), in radians: \code{th\_\allowbreak{}tarq} (angle between q and the
polarization,\allowbreak{} NaN when \code{targ\_\allowbreak{}pol\ =\ 0}), \code{phitarq}, \code{beta}, \code{phis} (Sivers angle φpq − φtarg), \code{phic} (Collins angle φpq

\begin{itemize}
\item φtarg), and the vertex values \code{betai}, \code{phisi}, \code{phici}. The kaon columns \code{saghai}/\code{factor} and the rho columns
\code{Mrho}/\code{Thrho}/\code{mmnuc} move after them. The shipped polarized-target pion example has 61 columns.
\end{itemize}
\subsection{\texorpdfstring{Values that are not measurements}{Values that are not measurements}}\label{user_manual-07u5foutputs:718-values-that-are-not-measurements}
\begin{itemize}
\item With an arm switched off (\hyperref[param:spectu5fmode]{\code{spect\_\allowbreak{}mode}}\simcidx{Parameters!spect_mode@\code{spect\_mode}} ≠ 0), that arm's focal-plane columns and \code{*ytar} are 0.0 in every event
(\hyperref[user_manual-06u5fspectrometers:64-spectu5fmode-switching-the-spectrometer-monte-carlo-off]{6.4}).
\item Other constant or NaN columns: \hyperref[user_manual-10u5ffaq:nan-or-constant-columns-in-the-tree]{10}.
\end{itemize}
\section{\texorpdfstring{\code{.\allowbreak{}hist}}{.hist}}\label{user_manual-07u5foutputs:72-hist}
The \code{.\allowbreak{}hist} is the run report, in the classic SIMC layout: the run title and
time stamps, the FLAGS block, the kinematics, the target, the central event and its cross section (\code{CENTRAL.\allowbreak{}sigcc}), the
radiative settings, the event counters (\code{Ngen\ (\allowbreak{}request)}, \code{Ntried}, \code{Ncontribute}, \code{Nco\_\allowbreak{}no\_\allowbreak{}rad\_\allowbreak{}prot}), the
normalisation (\code{wtcontr}, \code{AVERAGE.\allowbreak{}sigcc}, \code{charge}, \code{targetfac}, luminosity, \code{genvol}, \code{normfac}), the resolutions
(RECON SUMMARY), the input spectrometer limits (printed \textbf{after} the slop widening) and the limit
tables. Lines 2-3 are the BEGIN/END times; everything else depends only on the input, the seed and the mode.

\section{\texorpdfstring{\code{.\allowbreak{}geni}, \code{.\allowbreak{}gen} and \code{.\allowbreak{}gen.\allowbreak{}root}}{.geni, .gen and .gen.root}}\label{user_manual-07u5foutputs:73-geni-gen-and-genroot}
\begin{itemize}
\item \code{.\allowbreak{}geni}: for each spectrometer used, the number of trials and where they were stopped (slits, magnet entrances,
middles and exits, hut, chambers, scintillators,\allowbreak{} calorimeter),\allowbreak{} and the successes. Counter keys per spectrometer are
listed in \code{docs/\allowbreak{}manual/\allowbreak{}markdown/\allowbreak{}reference/\allowbreak{}OUTPUT\_\allowbreak{}FORMAT.\allowbreak{}md} §2. With \hyperref[param:usingu5fHMScoll]{\code{using\_\allowbreak{}HMScoll}}\simcidx{Parameters!using_HMScoll@\code{using\_HMScoll}}/\hyperref[param:usingu5fSHMScoll]{\code{using\_\allowbreak{}SHMScoll}}\simcidx{Parameters!using_SHMScoll@\code{using\_SHMScoll}} a \code{Collimator\ stops\ (\allowbreak{}pi/\allowbreak{}mu)} line is added.
\item \code{.\allowbreak{}gen}: diagnostic histograms (50 bins) of delta, yptar, xptar of both arms, Em and Pm: \code{geni} (all tries,
vertex), \code{gen} (accepted, vertex) and \code{recon} (accepted, reconstructed,\allowbreak{} weighted). \code{.\allowbreak{}gen.\allowbreak{}root} has the same histograms
as TH1D named \code{<set>\allowbreak{}\_\allowbreak{}<arm>\allowbreak{}\_\allowbreak{}<var>} (for example \code{gen\_\allowbreak{}e\_\allowbreak{}delta}, \code{recon\_\allowbreak{}Em}), filled with the \code{.\allowbreak{}gen} binning (no
under/\allowbreak{}overflow). \code{gen} and \code{geni} use −xptar and \code{recon} uses +xptar, and \code{gen\_\allowbreak{}Pm} is never filled.
\end{itemize}
\section{\texorpdfstring{Run summary (\code{.\allowbreak{}summary.\allowbreak{}json}, \code{.\allowbreak{}summary.\allowbreak{}txt})}{Run summary (.summary.\allowbreak{}json,\allowbreak{} .summary.\allowbreak{}txt)}}\label{user_manual-07u5foutputs:74-run-summary-summaryjson-summarytxt}
A JSON document (\code{"schema":\ "simc\_\allowbreak{}run\_\allowbreak{}summary"}, version 1, \code{docs/\allowbreak{}manual/\allowbreak{}markdown/\allowbreak{}reference/\allowbreak{}summary.\allowbreak{}schema.\allowbreak{}json}) with the run identification,\allowbreak{} the
counters, the normalisation,\allowbreak{} all FLAGS, the stop counters, the kinematics, the central event, the limits and the tree
location. \code{.\allowbreak{}summary.\allowbreak{}txt} is the same in readable form. the comparison tool \code{simc\_\allowbreak{}compare} reads it;
\code{hist2json} converts a \code{.\allowbreak{}hist}/\code{.\allowbreak{}geni} into the same format. A threaded run adds \code{"run":\ \{"threads":\ N\}}.

\section{\texorpdfstring{Exit status}{Exit status}}\label{user_manual-07u5foutputs:75-exit-status}
\begin{simctable}{small}{2}{l>{\raggedright\arraybackslash}p{\dimexpr 0.869\linewidth-2\tabcolsep\relax}}
%
\toprule\rowcolor{simchead}\simcth{status} & \simcth{meaning} \\\midrule\endfirsthead
\multicolumn{2}{@{}l}{\simccontinued}\\\toprule\rowcolor{simchead}\simcth{status} & \simcth{meaning} \\\midrule\endhead
\midrule\multicolumn{2}{r@{}}{\simccontinues}\\\endfoot
\bottomrule\endlastfoot
0 & success \\
1 & usage error (unknown option, missing input file, bad \code{-{}-{}threads}\simcidx{Command-line options!threads@\code{-{}-{}threads}} value) \\
2 & input or configuration error: unreadable file, rejected value (arm code 6, π0 from the neutron on H, rho on A ≥ 4, decay of a non-π/K hadron, \code{-{}-{}threads}\simcidx{Command-line options!threads@\code{-{}-{}threads}} with \code{-{}-{}fortran-{}compat}\simcidx{Command-line options!fortran-compat@\code{-{}-{}fortran-{}compat}} or RANLUX, ...) \\
3 & run error: a set-up check stopped the run (the message says \code{STOP\ during\ setup\ .\allowbreak{}.\allowbreak{}.}), a missing data file, an I/O error \\
\end{simctable}
